\documentclass[a4paper,11pt]{amsart}
\usepackage[margin=1in]{geometry}
\usepackage[T1]{fontenc}
\usepackage[utf8]{inputenc}
\usepackage{lmodern}
\usepackage{amsmath,amssymb}
\usepackage{booktabs}
\usepackage{longtable}
\usepackage{array}
\usepackage{hyperref}

\title{Report on Integrated Pull Requests}
\date{July 19, 2026}

\begin{document}
\maketitle

\begin{abstract}
This report summarizes the already integrated pull-request material in the
latest high-numbered integration band, roughly PRs \#775--\#985.  It is based
on the local integration ledger, committed experimental notes, formalization
packages, and representative source notes; it is not a fresh execution replay.
The main conclusion is that the PR stream is useful, but mostly as a proof
program: it narrows residual cases, repairs false statements, adds formalized
kernels, and improves audit discipline.  The most promotion-worthy material is
the Paper D v13.2 cap package, the Lean discharges for Paper D components, and
the KoalaBear M1 rank-nine route cuts.  The finite adjacent safe side and the
asymptotic entropy-frontier theorem remain open.
\end{abstract}

\section{Scope and Method}

The request was to look at PRs already integrated, not to merge a new open
queue.  I therefore used the local repository state, especially
\texttt{experimental/agents-log.md}, the recent commits ending at
\texttt{951ad203}, and the integrated notes under \texttt{experimental/notes/}.
The effective scope is the latest approximately 200-number band of already
handled PRs, from the L1/BC/L2 wave around \#775 through the M1 moving-root
boundary cut \#985.

No Python, Lean, Sage, or long-running certificate script was executed for this
report.  Status words such as PROVED, CONDITIONAL, COUNTEREXAMPLE, AUDIT, and
EXPERIMENTAL are inherited from the integrated notes and should be read subject
to their local hypotheses.  In particular, ``proved local'' means proved under
the packet's printed contracts; it does not by itself promote a main-paper
theorem or close the Proximity Prize threshold.

\section{Executive Assessment}

The integrated PRs are not noise.  They have turned the repository from a
collection of speculative ledgers into a much more structured proof program.
The best contributions fall into four groups.

First, the Paper D line has matured.  The v13.2 promotion gives a current
authority source for the cap paper, with identity-prefix unsafe edges and a
clear adjacent-safe certificate template.  It supersedes the old v12/raw-v13
split as the main Paper D target.

Second, the formalization stream is now genuinely useful.  Several Lean PRs do
not merely restate definitions: they refute false skeleton signatures, repair
interfaces, and discharge previously sorried kernels.  The strongest recent
example is the ECFFT constructive discharge: \texttt{prop\_rational\_floor} and
\texttt{cor\_ecfft\_onestep} are proved with byte-identical statements in the
Lean skeleton.  The phi-fiber repair certificate is also valuable because it
separates an old Lean-skeleton defect from the actual Paper D statement.

Third, the KoalaBear M1 branch is more concrete.  The sequence of rank-nine
route cuts now turns a broad full-outside high-degree residual into a small
list of named algebraic one-slack cases.  The latest cut, PR \#985, shows that
after the source-rational owner, every surviving full-outside rank-two record
has
\[
        s\ge 67{,}474,\qquad e\ge 33{,}737,\qquad
        \deg\gcd(P,Q)\le 1{,}014{,}838,
\]
and the remaining one-slack shapes are
\[
\begin{array}{c}
\texttt{UNPAID\_NONBASE\_COMMON\_LINEAR\_GCD\_TWIST},\\
\texttt{UNPAID\_SPLIT\_GCD\_NONBASE\_LINEAR\_MOVING\_COFACTOR}.
\end{array}
\]
This is useful because it tells the next proof exactly where to look.

Fourth, the counterexample and audit PRs are doing necessary cleanup.  They
prevent attractive but false shortcuts: wrong support universes, false strict
crossing demands, stale certificate bindings, and overbroad heavy-fiber claims.
Those results do not ``move the leaderboard'', but they make the proof program
less fragile.

\section{Integrated Waves}

\scriptsize
\begin{longtable}{>{\raggedright\arraybackslash}p{0.12\textwidth}
                  >{\raggedright\arraybackslash}p{0.28\textwidth}
                  >{\raggedright\arraybackslash}p{0.32\textwidth}
                  >{\raggedright\arraybackslash}p{0.10\textwidth}}
\toprule
PR range & Main content & Utility judgment & Promote now?\\
\midrule
\endhead
\#775, \#777, \#779--\#802 &
L1 mixed-petal and CRT packets, BC chart typing, resonant/fold/cylinder
reductions, L2 rank-15/rank-16 route cuts, and several Lean threshold
compilers. &
Useful for the final-resolution atlas.  It makes residual cells more
checkable, but it does not close Q/BC/SP or any deployed adjacent safe row. &
No; summarize selectively.\\
\addlinespace
\#803--\#815 &
L1 shared-auxiliary and total-degree CRT packets, rank-cell exclusions,
KoalaBear base-slope universe audits, BC deficiency-2 typing, and site updates
for corrected rows. &
High utility for finite ledgers and public-board accuracy.  The most visible
effect is better row display, not a new theorem closing the main gap. &
Site yes; paper selective.\\
\addlinespace
\#817--\#845 &
Threshold/compiler notes, C0 periodic Lean stack, M1 branch stack, L2
source certificates, fixed-slope Johnson formalization, and rank-one
corrections. &
Important infrastructure.  The Lean and compiler items are reusable, while
many route cuts remain local and conditional on packet contracts. &
Map only.\\
\addlinespace
\#849, \#867, \#871, \#874 &
KoalaBear branch-3 rank-nine route-cut stack: low-excess carrier cut, actual
core MDS rank ladder, mask-deficit route cut, and syndrome-rank reduction. &
High utility for M1.  These packets make the rank-nine branch structured, but
leave sparse-sigma/rank-nine residuals and the KoalaBear row open. &
No; keep as route cuts.\\
\addlinespace
\#879--\#880 &
Paving-v8 audit and dense-shell class-charge material. &
Useful for ePrint readiness and sign/class-sum separation.  It is supportive
for the paving paper rather than a direct final-threshold proof. &
Audit yes; theorem no.\\
\addlinespace
\#881--\#950 &
Large wave: Route-D no-go and recursion tools, dense-shell repairs,
transfer-shape evidence, M1 rank-nine route cuts, L2 rank-16 fixed-26/fixed-27
cuts, Paving retained-factor audits, and Lean progress. &
Mixed but useful.  The important effect is a stronger residual map and better
guardrails against false promotion.  Heavy computations should remain
off-by-default. &
No wholesale promotion.\\
\addlinespace
\#951--\#975 &
Route-D owner-typed forest and marked recursion, M1 source-Mobius and
projective-base refinements, L2 fixed-26/fixed-27 local packets, dense-shell
certificate audits, affine-line Steiner quotient ownership, and Lean repairs. &
Very useful locally.  PR \#962 is the strongest mathematical move in this
subwave, removing a selector-provenance gate for one M1 full-outside
maximal-gcd cell.  PR \#965 is important because it keeps stale Arb evidence in
OPEN-GAP status. &
Only audited summaries.\\
\addlinespace
\#976--\#979 &
Affine-prefix fixed-point and pole strata, primitive top-seam marked-incidence
audit, and complete-fiber subset-packing Lean package. &
Good regression material.  The top-seam audit is especially useful against
marked-incidence overcounting; the Lean subset-packing compiler is reusable.
No leaderboard movement. &
Lean maybe; theory no.\\
\addlinespace
\#980--\#984 &
Phi-fiber Lean repair certificate, ECFFT Lean constructive discharge, O5c
strict-crossing counterexample, Case-B support-universe correction, and M1
source-rational owner route cut. &
High value.  The Lean discharges and phi-fiber repair are audit/formalization
wins; the O5c and Case-B notes prevent wrong proof targets; the M1 owner cut
narrows the KoalaBear rank-nine residual. &
Audit.\\
\addlinespace
\#985 &
M1 moving-root cofactor/slack and C5 boundary cut. &
High local utility.  It raises the post-\#982 full-outside source floor to
\(s\ge 67{,}474\) and leaves two explicit one-slack residuals. &
Route-cut summary yes.\\
\bottomrule
\end{longtable}
\normalsize

\section{Most Useful Results}

\subsection{Paper D and threshold authority}

The main paper-level movement is not an individual small PR but the promotion
of \texttt{tex/cs25\_cap\_v13\_2.tex}.  It makes Paper D v13.2 the current
authority source.  Its useful contribution is the status split:
\[
\text{unsafe exact certificates are theorem-level in the paper package,}
\]
while adjacent first-safe statements remain conditional on exact upper ledgers.
For the deployed KoalaBear MCA row, the current experimental unsafe agreement
is \(a_0=1{,}116{,}047\), corresponding to
\[
        \delta=\frac{981105}{2097152}\approx 0.4678273.
\]
The next finite target is not to get a bigger error at a worse radius; it is to
prove the adjacent safe certificate at \(a_0+1=1{,}116{,}048\), with exact
integer constants.

\subsection{Formalization}

The integrated Lean material has crossed from outline to useful proof
infrastructure.  The most important formalization items are:
\begin{itemize}
\item PR \#961, which refutes false Paper-D-v12 skeleton statements, repairs
the formulations, and proves the \(c=1\) first-grid cap chain.
\item PR \#977, which formalizes the complete-fiber subset-packing compiler.
\item PR \#980, which gives a Lean counterexample to the old phi-fiber skeleton
signature and a coefficient-to-value wrapper matching the Paper D v13.2
premise \(\varphi\in B[X]\).
\item PR \#983, which proves \texttt{prop\_rational\_floor} and
\texttt{cor\_ecfft\_onestep} in the ECFFT skeleton, reducing the documented
sorry count.
\end{itemize}
These are among the most promotion-worthy PR outputs because they reduce
ambiguity rather than adding another heuristic layer.

\subsection{M1 KoalaBear safe-side route cuts}

The M1 packets are the strongest route toward the finite adjacent safe side.
Their value is cumulative.  Early branch-3 packets split the problem into
rank and carrier cases.  The source-Mobius/source-rational packets then remove
large projective-rational owner classes without increasing the paid bound.  The
moving-root C5 boundary packet composes the cofactor divisibility with the
source-rational survivor inequality.

The important current local endpoint is:
\[
        \text{post-C5 full-outside rank-two survivors have }
        s\ge 67{,}474,\ e\ge 33{,}737.
\]
This replaces a vague high-degree residual by two explicit one-slack
components.  The next useful theorem is therefore sharply stated: classify
those two nonbase linear components against the regular locator equations and
the eight rank-nine outlier directions.  Generic high-degree elimination is not
the next best move.

\subsection{L1, BC, L2, and Route-D}

The L1/BC/L2/Route-D material is broad.  Its most useful function is not a
single theorem but a better residual atlas.  The PRs add mixed-petal and CRT
evidence on the L1 side, BC chart typing and first-interior modular-fiber
notes, fixed-slope and complete-fiber L2 compilers, and Route-D owner/no-go
packets.

The judgment is conservative: these should not be promoted wholesale into the
main TeX papers.  They are best used as a menu of local lemmas and
counterexamples.  A future paper update should extract only statements that
either close a named residual branch or give a reusable theorem with clearly
printed hypotheses.

\subsection{Audits and counterexamples}

Several integrated PRs are valuable because they say ``do not prove this'':
\begin{itemize}
\item PR \#981 refutes a universal O5c strict-crossing demand using a sparse
exact-deep family where the exact numerator is \(2\), equal to or below the
frozen target in the relevant endpoint.
\item PR \#984 corrects a Case-B support-universe mistake: supports must be
counted by \(\binom{|D|}{m}\), not by \(\binom{|B|}{m}\).  The corrected
fixed-density regime still gives an exponential-field lower picture, but the
old \(m=2e\) calibration was wrong.
\item PR \#978 is a useful primitive top-seam regression against marked
incidence overcounting and common-core erasure.
\item PR \#965 is useful precisely because it refuses to promote stale Arb
evidence to a proved transfer theorem.
\end{itemize}
These negative or audit results should stay visible.  They prevent future
agents from spending cycles on false shortcuts.

\section{What Should Be Promoted}

The following items are strong candidates for promotion into a clean TeX
writeup, after human review.

\begin{enumerate}
\item The ECFFT Lean discharge can be summarized in the Paper D formalization
appendix or an audit note: the rational floor and one-step ECFFT cap have
machine-checked skeleton proofs with byte-identical statements.
\item The phi-fiber repair certificate should be cited in the Paper D audit
map as a formal check that the base-field polynomial premise is load-bearing
and correctly bridges to the repaired value-level Lean theorem.
\item The M1 KoalaBear route-cut sequence, ending with PR \#985, should be
summarized as the current best finite safe-side residual reduction.  It should
not be advertised as KoalaBear closure.
\item The Case-B correction should be summarized in the asymptotic/frontier
notes because it changes the correct support-counting calibration and
separates subexponential-target lower constructions from fixed-epsilon
constant-fraction questions.
\item The complete-fiber subset-packing Lean compiler from PR \#977 should be
kept in the formalization roadmap as reusable L2 infrastructure.
\end{enumerate}

\section{What Should Not Be Promoted Yet}

The following should remain experimental.

\begin{enumerate}
\item Local route cuts whose hypotheses depend on selector provenance,
first-match owner order, or frozen packet contracts.
\item Large Python/Sage certificate packets unless a maintainer explicitly
requests replay and the certificate is source-bound.
\item Leaderboard or site rows that are only evidence, not theorem-level
unsafe/safe certificates.
\item Broad L1/BC/L2/Route-D stacks that narrow residuals but do not close a
named theorem.
\item Any statement using unspecified polynomial losses to close a finite
adjacent deployed row.  The finite rows need constants fitting inside the
printed bit margin.
\end{enumerate}

\section{Current Research State After These PRs}

For the finite deployed Proximity Prize direction, the unsafe side is much
sharper than before.  The main remaining task is the adjacent safe certificate:
if \(B^*=\lfloor 2^{-128}q_{\rm line}\rfloor\), prove
\[
        U(a_0+1)\le B^* < L(a_0)
\]
with an exhaustive, deduplicated upper ledger \(U\) and an exact lower
staircase \(L\).  The PRs mostly improve \(U\), especially through M1 route
cuts.

For the asymptotic direction, the conjectural picture is still the
entropy-subfield frontier
\[
        \delta_C^*(\varepsilon^*)
        =
        1-\rho-g^*(\rho,\log_2 |B|)+o(1),
\]
where \(g^*\) is determined by the binary-entropy balance.  The recent PRs do
not prove this theorem, but they improve the evidence and remove several false
formulations of the required support-counting and residual-payment inputs.

For formalization, the project is clearly moving in the right direction.  The
next best formalization work is not to add more skeletons, but to mechanize
the exact staircase compilers, the Paper D v13.2 cap chain, and the M1
route-cut statements with their dependencies explicitly printed.

\section{Bottom Line}

The last integrated PR band is useful and worth keeping, but most of it should
remain in \texttt{experimental/}.  The material most likely to survive into a
submission-facing paper is:
\[
\begin{aligned}
&\text{Paper D v13.2 authority}
+\text{ Lean formalization discharges}\\
&\quad+\text{ KoalaBear M1 residual route cuts}
+\text{ support-universe and strict-crossing audits}.
\end{aligned}
\]
The project is closer to RS MCA because the remaining proof obligations are
now narrower and better named.  It is not done: the finite adjacent safe
certificate and the asymptotic entropy-frontier upper theorem still require
new proof, not just more certificates.

\end{document}
